% =====================================================================
%  CARNET D'ENTRAINEMENT — FICHE 6 — THÈME 3
%  Niveau DIFFICILE — Corrigé
% =====================================================================
\documentclass[11pt,a4paper]{article}
\usepackage[utf8]{inputenc}
\usepackage[T1]{fontenc}
\usepackage{lmodern}
\usepackage[french]{babel}
\usepackage{amsmath,amssymb,mathtools}
\usepackage{icomma}
\usepackage[version=4]{mhchem}
\usepackage{siunitx}
\sisetup{output-decimal-marker={,}}
\usepackage{xcolor}
\usepackage{colortbl}
\usepackage{tikz}
\usepackage[most]{tcolorbox}
\usepackage{enumitem}
\usepackage{array}
\usepackage{booktabs}
\usepackage{fancyhdr}
\usepackage[hidelinks]{hyperref}
\usetikzlibrary{arrows.meta,positioning,calc,shapes.geometric}
\usepackage[a4paper,margin=2.1cm,headheight=26pt,headsep=8pt]{geometry}

\definecolor{primary}{RGB}{146,95,160}
\definecolor{primarydark}{RGB}{92,52,104}
\definecolor{corcol}{RGB}{0,114,178}
\definecolor{astucecol}{RGB}{198,93,9}

\tcbset{boxbase/.style={enhanced,breakable,boxrule=0.9pt,arc=2.5pt,
   left=3mm,right=3mm,top=2.2mm,bottom=2.2mm,
   fonttitle=\bfseries,coltitle=white,toptitle=0.6mm,bottomtitle=0.6mm}}
\newtcolorbox[auto counter]{cor}[1][]{boxbase,colback=corcol!5,colframe=corcol,
   title={\textsc{Corrigé}~\thetcbcounter\ifx\relax#1\relax\relax\else\textmd{ —\itshape\ #1}\fi}}
\newtcolorbox{astucebox}[1][]{boxbase,colback=astucecol!8,colframe=astucecol,
   title={\textsc{Astuce}\ifx\relax#1\relax\relax\else\textmd{ : \itshape #1}\fi}}

\pagestyle{fancy}\fancyhf{}
\fancyhead[L]{\small\textcolor{primarydark}{\textbf{Fiche 6 · Thème 3} — Corrigé}}
\fancyhead[R]{\small\textcolor{primarydark}{\textsc{Difficile}}}
\fancyfoot[C]{\small\thepage}
\renewcommand{\headrulewidth}{0.8pt}
\renewcommand{\headrule}{\hbox to\headwidth{\color{primary}\leaders\hrule height \headrulewidth\hfill}}
\setlength{\parindent}{0pt}\setlength{\parskip}{3pt}

\newcommand{\rep}{\textbf{\textcolor{corcol}{Réponse : }}}

\begin{document}
\begin{center}
{\Large\bfseries\color{primarydark}Thème 3 — Corrigé du niveau Difficile}
\end{center}
\medskip

\begin{cor}[ion perchlorate \ce{ClO4-} : deux structures concurrentes]
\begin{enumerate}[label=\alph*),leftmargin=2em,itemsep=3pt]
  \item \textbf{Structure A} (4 simples liaisons).\\
        Chlore : $0$ doublet libre, $4$ liaisons $\Rightarrow
        q_{\mathrm{f}}(\ce{Cl}) = 7 - 2\times 0 - 4 = \mathbf{+3}$.\\
        Chaque \ce{O} simplement lié : $3$ doublets libres, $1$ liaison $\Rightarrow
        q_{\mathrm{f}}(\ce{O}) = 6 - 2\times 3 - 1 = \mathbf{-1}$.\\
        Somme $= (+3) + 4\times(-1) = 3 - 4 = \mathbf{-1}$. \checkmark
  \item \textbf{Structure B} (octet étendu : 3 doubles $+$ 1 simple).\\
        Chlore : $0$ doublet libre, $7$ liaisons (3 doubles $=6$, $+1$ simple) $\Rightarrow
        q_{\mathrm{f}}(\ce{Cl}) = 7 - 2\times 0 - 7 = \mathbf{0}$.\\
        Chaque \ce{O} doublement lié : $2$ doublets libres, $2$ liaisons $\Rightarrow
        q_{\mathrm{f}} = 6 - 2\times 2 - 2 = \mathbf{0}$.\\
        L'\ce{O} simplement lié : $3$ doublets libres, $1$ liaison $\Rightarrow
        q_{\mathrm{f}} = 6 - 2\times 3 - 1 = \mathbf{-1}$.\\
        Somme $= 0 + 3\times 0 + (-1) = \mathbf{-1}$. \checkmark
  \item \rep La structure \textbf{B} est la plus représentative. Critère : à octet respecté (ou
        étendu par un atome qui en a le droit), on retient la structure aux \textbf{charges
        formelles minimales}. La structure B ramène le chlore de $+3$ à $0$ et annule les charges de
        trois oxygènes : ses charges formelles sont bien plus proches de zéro que celles de A.
  \item \rep Passer de A à B, c'est ajouter $3$ liaisons $\pi$ au chlore, qui s'entoure alors de
        $14$ électrons ($7$ liaisons $\times 2$) : c'est un \textbf{octet étendu}. Seuls les
        éléments de \textbf{période $\geq 3$} (comme \ce{Cl}, qui dispose d'orbitales $d$
        accessibles) peuvent dépasser $8$ électrons. La structure B serait \emph{interdite} pour un
        atome central de 2\textsuperscript{e} période.
\end{enumerate}
\end{cor}

\begin{cor}[choisir la bonne structure : l'acide nitrique \ce{HNO3}]
\begin{enumerate}[label=\alph*),leftmargin=2em,itemsep=3pt]
  \item L'azote forme \textbf{une double} \ce{N=O} $+$ \textbf{deux simples} \ce{N-O}, soit
        $2 + 1 + 1 = \mathbf{4}$ liaisons et aucun doublet libre. Il s'entoure de
        $4\times 2 = \mathbf{8}$ électrons : \rep son octet est \emph{exactement} respecté.
  \item \rep L'azote est en \textbf{2\textsuperscript{e} période} : il ne possède pas d'orbitales
        $d$ et ne peut donc \textbf{jamais} dépasser $8$ électrons. Une structure à $5$ liaisons lui
        ferait porter $10$ électrons — impossible. Toute structure « à azote pentavalent » est
        à rejeter (critère 1 : octet obligatoire pour la 2\textsuperscript{e} période).
  \item Charges formelles dans la bonne structure :\\
        Azote : $0$ doublet libre, $4$ liaisons $\Rightarrow
        q_{\mathrm{f}}(\ce{N}) = 5 - 2\times 0 - 4 = \mathbf{+1}$.\\
        \ce{O} doublement lié : $2$ doublets libres, $2$ liaisons $\Rightarrow q_{\mathrm{f}} = 6-4-2 = 0$.\\
        \ce{O} du groupe \ce{O-H} : $2$ doublets libres, $2$ liaisons (\ce{N-O} et \ce{O-H})
        $\Rightarrow q_{\mathrm{f}} = 6-4-2 = 0$.\\
        \ce{O} simplement lié \emph{sans} \ce{H} : $3$ doublets libres, $1$ liaison $\Rightarrow
        q_{\mathrm{f}} = 6 - 2\times 3 - 1 = \mathbf{-1}$.\\
        Chaque \ce{H} : $q_{\mathrm{f}} = 0$.\\
        \rep Somme $= (+1) + 0 + 0 + (-1) + 0 = \mathbf{0}$. \checkmark\; Cohérent avec une molécule
        neutre : la charge $+$ (sur \ce{N}) et la charge $-$ (sur l'\ce{O} terminal) se compensent.
\end{enumerate}
\end{cor}

\begin{cor}[ion nitrate \ce{NO3^-} : pourquoi l'azote garde-t-il sa charge ?]
\begin{enumerate}[label=\alph*),leftmargin=2em,itemsep=3pt]
  \item \rep Azote central, \textbf{une double} \ce{N=O} $+$ \textbf{deux simples} \ce{N-O^-},
        sans doublet libre : l'azote forme $2 + 1 + 1 = \mathbf{4}$ liaisons (soit $8$ électrons,
        octet strict).
  \item Azote : $q_{\mathrm{f}}(\ce{N}) = 5 - 2\times 0 - 4 = \mathbf{+1}$.\\
        \ce{O} doublement lié : $q_{\mathrm{f}} = 6 - 2\times 2 - 2 = \mathbf{0}$.\\
        Chaque \ce{O} simplement lié : $q_{\mathrm{f}} = 6 - 2\times 3 - 1 = \mathbf{-1}$.\\
        \rep Somme $= (+1) + 0 + 2\times(-1) = 1 - 2 = \mathbf{-1}$. \checkmark\; C'est bien la charge
        de l'ion.
  \item \rep La structure est « bonne » parce que l'azote, en 2\textsuperscript{e} période,
        \textbf{respecte strictement l'octet} ($4$ liaisons $= 8$ électrons) : c'est le maximum
        permis. Contrairement au chlore de \ce{ClO4-} (3\textsuperscript{e} période), l'azote
        \textbf{ne peut pas} former de doubles liaisons supplémentaires pour annuler sa charge
        $+1$ : cela dépasserait l'octet, ce qui lui est interdit. Le $+1$ sur l'azote est donc
        \emph{inévitable} et parfaitement acceptable ici.
\end{enumerate}
\end{cor}

\begin{astucebox}[le réflexe de la charge formelle]
$q_{\mathrm{f}} = N_v - 2\times(\text{doublets non liants}) - (\text{nb de liaisons})$, et
\textbf{une double liaison compte pour 2 liaisons}. L'octet étendu (pour abaisser une charge) n'est
permis qu'aux atomes de \textbf{période $\geq 3$}. Vérifie toujours : la somme des $q_{\mathrm{f}}$
doit redonner la charge de l'édifice.
\end{astucebox}

\end{document}
